Volem utilitzar un espectròmetre de masses per a separar els ions de carboni $^{12}_{6}\mathrm{C}^-$ i $^{14}_{6}\mathrm{C}^-$, amb la finalitat de datar un fòssil. En un gresol, sublimem el material i posteriorment, emprant un feix d'electrons, els àtoms són ionitzats, de manera que cada ió ha adquirit una càrrega $-e$. A continuació, fem passar aquests ions per un selector de velocitats, amb la qual cosa tots els ions entren dins l'espectròmetre amb la mateixa velocitat. Finalment, els ions són desviats per un camp magnètic, de manera que cada tipus d'ió surt per una obertura diferent (a la imatge següent es veu una d'aquestes obertures).

\figura[0.4]{fig1.png}

\begin{apartat}{1,25}
A partir de la força magnètica, determineu el radi de la trajectòria en funció de la massa de l'ió, $m$, de la seva velocitat, $v$, de la seva càrrega, $e$, i del mòdul del camp magnètic, $B$; és a dir, deduïu com s'expressa el radi de la trajectòria en funció de $m$, $v$, $e$ i $B$.
\end{apartat}

\begin{apartat}{1,25}
Sabem que tots els ions entren dins l'espectròmetre amb una velocitat de $4,80\times 10^5\ \mathrm{m/s}$ i que els ions $^{12}_{6}\mathrm{C}^-$ surten per una obertura situada a 30,0 cm de l'entrada. Calculeu el mòdul del camp magnètic. A quina distància de l'entrada hem de posar la segona obertura per a recollir els ions $^{14}_{6}\mathrm{C}^-$?
\end{apartat}

\nota{En el càlcul de la massa dels ions, podeu negligir la massa dels electrons i suposar que els protons i els neutrons tenen la mateixa massa.}

\dades{%
  $m_\mathrm{p} = 1,67\times 10^{-27}\ \mathrm{kg}$.\\
  $|e| = 1,602\times 10^{-19}\ \mathrm{C}$.}
